\section*{LISTA DE SIGLAS Y ABREVIATURAS}

{\renewcommand{\arraystretch}{1.15}
\begin{longtable}{@{}>{\bfseries}p{2.4cm}p{12.6cm}@{}}
3DS & Protocolo con el que el banco emisor autentica al titular de la tarjeta en una compra en línea (\textit{Three-Domain Secure}). \\
ACH & Cámara de compensación automatizada (\textit{Automated Clearing House}). En este documento, ACH Colombia, la entidad que opera PSE. \\
ADR & Registro de una decisión de arquitectura (\textit{Architecture Decision Record}). \\
APA & Asociación Estadounidense de Psicología (\textit{American Psychological Association}), cuyo estilo de citas y referencias sigue este documento. \\
API & Interfaz de programación de aplicaciones (\textit{Application Programming Interface}). \\
ASR & Requisito arquitectónicamente significativo (\textit{Architecturally Significant Requirement}). \\
BIN & Número de identificación bancaria, formado por los primeros dígitos de la tarjeta (\textit{Bank Identification Number}). \\
C4 & Modelo que describe una arquitectura en cuatro niveles: contexto, contenedores, componentes y código (\textit{Context, Containers, Components, Code}). \\
CBO & Acoplamiento entre objetos (\textit{Coupling Between Objects}). \\
CC & Complejidad ciclomática de una función. \\
CI & Integración continua (\textit{Continuous Integration}). \\
CK & Conjunto de métricas de diseño orientado a objetos de Chidamber y Kemerer (1994). \\
COP & Código ISO 4217 del peso colombiano. \\
CORS & Intercambio de recursos de origen cruzado (\textit{Cross-Origin Resource Sharing}). \\
CP & Caso de prueba del plan de pruebas (CP-01 a CP-07). \\
CVV & Código de verificación de la tarjeta (\textit{Card Verification Value}). \\
DANE & Departamento Administrativo Nacional de Estadística. \\
DDD & Diseño guiado por el dominio (\textit{Domain-Driven Design}). \\
DEF & Defecto registrado en la ejecución del plan de pruebas (DEF-01 y DEF-02). \\
DIT & Profundidad del árbol de herencia de una clase (\textit{Depth of Inheritance Tree}). \\
DoD & Definición de terminado (\textit{Definition of Done}). \\
DSR & Investigación en ciencia del diseño (\textit{Design Science Research}). \\
ERP & Planificación de recursos empresariales (\textit{Enterprise Resource Planning}). \\
et al. & Y otros, del latín \textit{et alii}, en las citas con más de dos autores. \\
FURPS+ & Modelo de clasificación de requisitos no funcionales: funcionalidad, usabilidad, confiabilidad, rendimiento y soporte, más las restricciones (\textit{Functionality, Usability, Reliability, Performance, Supportability}). \\
GoF & Banda de los Cuatro, nombre con que se conoce a los autores del catálogo de patrones de diseño de Gamma et al. (1994) (\textit{Gang of Four}). \\
GPO & Organización global de pagos, la unidad de PayU que compró Rapyd (\textit{Global Payment Organisation}). \\
H1 a H6 & Hitos del ciclo de vida del proyecto. \\
HMAC & Código de autenticación de mensajes basado en una función de resumen (\textit{Hash-based Message Authentication Code}). \\
HTML & Lenguaje de marcado de hipertexto (\textit{HyperText Markup Language}). \\
HTTP & Protocolo de transferencia de hipertexto (\textit{Hypertext Transfer Protocol}). \\
HTTPS & HTTP sobre una conexión cifrada con TLS (\textit{Hypertext Transfer Protocol Secure}). \\
HU & Historia de usuario del SRS (HU-01 a HU-13). \\
IEEE & Instituto de Ingenieros Eléctricos y Electrónicos (\textit{Institute of Electrical and Electronics Engineers}). \\
IP & Protocolo de internet (\textit{Internet Protocol}). En este documento, la dirección IP del pagador. \\
ISO & Organización Internacional de Normalización (\textit{International Organization for Standardization}). \\
IVA & Impuesto al valor agregado. \\
JSON & Notación de objetos de JavaScript, un formato de intercambio de datos (\textit{JavaScript Object Notation}). \\
KB & Kilobyte. \\
LCOM & Falta de cohesión entre los métodos de una clase (\textit{Lack of Cohesion in Methods}). \\
LOC & Líneas de código (\textit{Lines of Code}). \\
MAX\_CC & Complejidad ciclomática más alta entre las funciones de una clase o de un archivo. \\
MIT & Instituto Tecnológico de Massachusetts, que da nombre a una licencia de código abierto (\textit{Massachusetts Institute of Technology}). \\
NASA & Administración Nacional de Aeronáutica y del Espacio de Estados Unidos (\textit{National Aeronautics and Space Administration}). \\
NOC & Número de hijos, o subclases directas, de una clase (\textit{Number of Children}). \\
OE & Objetivo específico del trabajo de grado (OE-01 a OE-06). \\
PAN & Número principal de la tarjeta (\textit{Primary Account Number}). \\
PCI DSS & Estándar de seguridad de datos de la industria de tarjetas de pago (\textit{Payment Card Industry Data Security Standard}). \\
pp & Puntos porcentuales. \\
PSE & Pagos Seguros en Línea, el servicio de ACH Colombia con el que el comprador paga con un débito a su cuenta bancaria. \\
PTP & Plan de Pruebas del Proyecto. \\
QR & Código de respuesta rápida (\textit{Quick Response}). \\
rc & Versión candidata a publicación (\textit{release candidate}). \\
REST & Transferencia de estado representacional, un estilo de diseño de servicios web (\textit{Representational State Transfer}). \\
RF & Requisito funcional del SRS (RF-01 a RF-19). \\
RFC & Respuesta de una clase (\textit{Response For a Class}). \\
SAD & Documento de arquitectura de software (\textit{Software Architecture Document}). \\
SAQ & Cuestionario de autoevaluación de PCI DSS (\textit{Self-Assessment Questionnaire}). \\
SATC & Centro de Aseguramiento de Tecnología de Software de la NASA (\textit{Software Assurance Technology Center}). \\
SDK & Kit de desarrollo de software (\textit{Software Development Kit}). \\
s.~f. & Sin fecha, en las citas y referencias. \\
SHA & Algoritmo seguro de resumen criptográfico (\textit{Secure Hash Algorithm}). \\
SPMP & Plan de gestión del proyecto de software (\textit{Software Project Management Plan}). \\
SRS & Especificación de requisitos de software (\textit{Software Requirements Specification}). \\
TLS & Seguridad de la capa de transporte (\textit{Transport Layer Security}). \\
UAT & Pruebas de aceptación de usuario (\textit{User Acceptance Testing}). En Kushki, el nombre de su ambiente de pruebas. \\
URL & Localizador uniforme de recursos, la dirección de una página o de un servicio web (\textit{Uniform Resource Locator}). \\
UUID & Identificador único universal (\textit{Universally Unique Identifier}). \\
WMC & Métodos ponderados por clase (\textit{Weighted Methods per Class}). \\
\end{longtable}}
\addtocounter{table}{-1}
